\topic{悬空输入为什么会变成随机按键}

输入引脚悬空时仍有微小电容，也会通过漏电、手指和邻近信号获得电荷。输入
缓冲器看到一个不受保证的模拟电压；当它穿过阈值时，软件可能观察到重复中断、
随机按键或复位抖动。

\begin{center}
\begin{tikzpicture}[font=\footnotesize\sffamily]
  \node[os hardware] (vcc) {\(\SI{3.3}{V}\)};
  \node[draw=BookAmber,minimum width=18mm,minimum height=7mm,right=14mm of vcc] (r) {\(R_{pull}\)};
  \node[circle,fill=BookBlue,inner sep=2pt,right=14mm of r] (point) {};
  \node[os hardware,right=14mm of point] (input) {CPU 输入引脚};
  \node[os hardware,below=14mm of point] (button) {按键接 GND};
  \draw[BookBlue,thick] (vcc) -- (r) -- (point) -- (input);
  \draw[BookBlue,thick] (point) -- (button);
\end{tikzpicture}
\end{center}

\begin{workedexample}
\(\SI{3.3}{V}\) 电源配 \(\SI{10}{k\ohm}\) 上拉。按键闭合时：
\[
I=\frac{V}{R}=\frac{\SI{3.3}{V}}{\SI{10}{k\ohm}}
=\SI{0.33}{mA},
\qquad
P=VI\approx\SI{1.09}{mW}.
\]
若误用 \(\SI{10}{\ohm}\)，电流变为 \(\SI{330}{mA}\)，功耗约
\(\SI{1.09}{W}\)。单位错三个数量级会把逻辑问题变成发热问题。
\end{workedexample}

\begin{workedexample}
输入漏电最坏为 \(\SI{5}{\micro A}\)，上拉为
\(\SI{100}{k\ohm}\)。电阻压降：
\[
\Delta V=IR=\SI{5}{\micro A}\times\SI{100}{k\ohm}
=\SI{0.5}{V}.
\]
节点最低约 \(\SI{2.8}{V}\)。若 \(V_{IH(min)}=\SI{2.0}{V}\)，静态裕量
为 \(\SI{0.8}{V}\)。还要检查输入电容带来的上升时间；电阻越大，静态电流
越小，但恢复越慢。
\end{workedexample}

\topic{RC 复位要算跨过阈值的时刻}

\[
V(t)=V_{CC}\left(1-e^{-t/(RC)}\right).
\]
复位监控器在输入达到 \(0.7V_{CC}\) 时释放：
\[
0.7=1-e^{-t/(RC)}
\Rightarrow t=-RC\ln(0.3)\approx1.204RC.
\]
若 \(R=\SI{10}{k\ohm}\)、\(C=\SI{1}{\micro F}\)，则
\(RC=\SI{10}{ms}\)，释放时刻约 \(\SI{12.04}{ms}\)。

\begin{center}
\begin{tikzpicture}[x=1.05cm,y=1.0cm,font=\footnotesize\sffamily]
  \draw[->,BookMuted] (0,0) -- (7.2,0) node[right]{时间};
  \draw[->,BookMuted] (0,0) -- (0,4.1) node[above]{电压};
  \draw[domain=0:6.7,smooth,variable=\x,BookBlue,very thick]
    plot ({\x},{3.5*(1-exp(-\x/1.6))});
  \draw[dashed,BookRed] (0,2.45) -- (6.8,2.45);
  \draw[dashed,BookRed] (1.93,0) -- (1.93,2.45);
  \node[anchor=east] at (0,2.45) {\(0.7V_{CC}\)};
  \node[below] at (1.93,0) {\(1.204RC\)};
\end{tikzpicture}
\end{center}

RC 解释延迟来源，但不保证慢上电、掉电重启和噪声下都有干净边沿。真实板卡常
使用带迟滞和电源良好输入的复位监控器。QEMU 从定义好的处理器复位状态开始，
不模拟这条充电曲线。

\topic{全加器、MUX 和译码器怎样进入 CPU}

一位全加器满足：
\[
S=A\oplus B\oplus C_{in},\qquad
C_{out}=AB+C_{in}(A\oplus B).
\]

\begin{center}
\begin{osgridtabular}{ccccc}
\(A\) & \(B\) & \(C_{in}\) & \(S\) & \(C_{out}\) \\ \hline
0 & 0 & 0 & 0 & 0 \\ \hline
0 & 0 & 1 & 1 & 0 \\ \hline
0 & 1 & 0 & 1 & 0 \\ \hline
0 & 1 & 1 & 0 & 1 \\ \hline
1 & 0 & 0 & 1 & 0 \\ \hline
1 & 0 & 1 & 0 & 1 \\ \hline
1 & 1 & 0 & 0 & 1 \\ \hline
1 & 1 & 1 & 1 & 1 \\ \hline
\end{osgridtabular}
\end{center}

\begin{workedexample}
用四个全加器计算
\(\texttt{1011}_2+\texttt{0110}_2\)：

\begin{center}
\begin{osgridtabular}{ccccc}
bit & \(A_i\) & \(B_i\) & \(C_{in}\) & \(S_i,C_{out}\) \\ \hline
0 & 1 & 0 & 0 & 1,0 \\ \hline
1 & 1 & 1 & 0 & 0,1 \\ \hline
2 & 0 & 1 & 1 & 0,1 \\ \hline
3 & 1 & 0 & 1 & 0,1 \\ \hline
\end{osgridtabular}
\end{center}

低四位为 \texttt{0001}，最外层进位为 1，五位和为
\texttt{10001}\(_2=17\)。四位寄存器写回 \texttt{0001}，CF 保存进位。
\end{workedexample}

MUX 决定下一状态来自顺序地址还是跳转目标：
\[
PC_{next}=
\begin{cases}
target,& branch\_taken=1,\\
PC+length,& branch\_taken=0.
\end{cases}
\]
译码器把操作码变成互斥控制信号。操作码 \texttt{41} 表示 JZ 时，应激活
PC 选择，不得同时激活 RAM 写使能。

\begin{pdfdiagram}
  \node[os state] (opcode) {操作码 \texttt{41}};
  \node[os block,right=17mm of opcode] (decoder) {译码器};
  \node[os state,right=17mm of decoder,yshift=15mm] (alu) {ALU=保持};
  \node[os state,right=17mm of decoder] (pc) {PC MUX=Z};
  \node[os state,right=17mm of decoder,yshift=-15mm] (memory) {RAM write=0};
  \draw[os arrow] (opcode) -- (decoder);
  \draw[os arrow] (decoder) -- (alu);
  \draw[os arrow] (decoder) -- (pc);
  \draw[os arrow] (decoder) -- (memory);
\end{pdfdiagram}

\topic{D 触发器保存的是边沿附近的输入}

\begin{center}
\begin{tikzpicture}[x=0.85cm,y=0.68cm,font=\footnotesize\sffamily]
  \foreach \y/\name in {3/CLK,2/D,1/Q} {
    \node[anchor=east] at (-0.3,\y) {\name};
    \draw[BookRule] (0,\y-0.35) -- (10,\y-0.35);
  }
  \draw[BookBlue,thick]
    (0,3) -- (2,3) -- (2,3.35) -- (3,3.35) -- (3,3)
    -- (6,3) -- (6,3.35) -- (7,3.35) -- (7,3) -- (10,3);
  \draw[BookTeal,thick]
    (0,2) -- (1,2) -- (1,2.35) -- (4.8,2.35) -- (4.8,2) -- (10,2);
  \draw[BookAmber,thick]
    (0,1) -- (2.1,1) -- (2.1,1.35) -- (6.1,1.35) -- (6.1,1) -- (10,1);
  \draw[dashed,BookRed] (2,0.5) -- (2,3.8);
  \draw[dashed,BookRed] (6,0.5) -- (6,3.8);
  \node[below] at (2,0.5) {采样 1};
  \node[below] at (6,0.5) {采样 0};
\end{tikzpicture}
\end{center}

Q 稍晚于边沿变化，这段时间是 clock-to-Q。D 必须在边沿前满足 setup，在边沿
后满足 hold。若边沿在 \(\SI{20}{ns}\)，
\(t_{setup}=\SI{1.2}{ns}\)，\(t_{hold}=\SI{0.5}{ns}\)，D 必须从
\(\SI{18.8}{ns}\) 到 \(\SI{20.5}{ns}\) 保持稳定。D 在
\(\SI{19.4}{ns}\) 才变化就违反 setup，内部节点可能进入亚稳态。

16 个 D 触发器共享时钟和写使能可组成 16-bit 寄存器。PC 在 D 输入前再放一
个 MUX，选择复位值、顺序地址、跳转目标或返回地址。教学 CPU 的正式 C++20
例子把修改集中到写回阶段：

\lstinputlisting[
  language=C++,
  firstline=252,
  lastline=285,
  caption={JZ、CALL 和 RET 最后才更新 PC}
]{../../examples/teaching_cpu/source/teaching_cpu.cpp}

\topic{连续动手：从加法器搭到四位计数器}

这组练习只增加一层硬件再观察一次状态，不允许直接把“计数器”当成现成模块。
可以在纸上画门，也可以使用任意数字逻辑仿真器；无论用什么工具，都要保存每
一级的真值表和时钟波形。

\begin{center}
\begin{tikzpicture}[font=\footnotesize\sffamily,node distance=8mm and 9mm]
  \node[os state] (half) {半加器};
  \node[os block,right=of half] (full) {全加器};
  \node[os block,right=of full] (adder) {四位加法器};
  \node[os state,below=of adder] (sr) {SR 锁存器};
  \node[os block,left=of sr] (dlatch) {D latch};
  \node[os block,left=of dlatch] (dff) {D 触发器};
  \node[os state,below=of dff] (register) {四位寄存器};
  \node[os hardware,right=of register] (counter) {四位计数器};
  \draw[os arrow] (half) -- (full);
  \draw[os arrow] (full) -- (adder);
  \draw[os arrow] (adder) -- (sr);
  \draw[os arrow] (sr) -- (dlatch);
  \draw[os arrow] (dlatch) -- (dff);
  \draw[os arrow] (dff) -- (register);
  \draw[os arrow] (register) -- (counter);
\end{tikzpicture}
\end{center}
\diagramnote{上半条链产生“下一值”，下半条链保存“当前值”。计数器只是把
寄存器当前值送回加法器，并在下一个边沿写回。}

\mustunderstand\textbf{第一步：先让一位加法器正确。}
半加器给出
\[
S=A\oplus B,\qquad C=AB.
\]
全加器可以由两个半加器组成：第一只计算 \(A+B\)，第二只把
\(C_{in}\) 加到第一次的和上，两个进位再做 OR。逐行核对前面的八行真值表，
尤其要测 \((A,B,C_{in})=(1,1,1)\)，它应得到 \(S=1,C_{out}=1\)。

\mustunderstand\textbf{第二步：串接四个全加器。}
最低位的 \(C_{in}=0\)，每一位的 \(C_{out}\) 接到更高一位的
\(C_{in}\)。手算 \(\texttt{0111}+\texttt{0001}\)：

\begin{osgridtabular}{L{0.08\linewidth}L{0.10\linewidth}L{0.10\linewidth}L{0.12\linewidth}L{0.12\linewidth}L{0.15\linewidth}}
bit & \(A_i\) & \(B_i\) & \(C_{in}\) & \(S_i\) & \(C_{out}\) \\ \hline
0 & 1 & 1 & 0 & 0 & 1 \\ \hline
1 & 1 & 0 & 1 & 0 & 1 \\ \hline
2 & 1 & 0 & 1 & 0 & 1 \\ \hline
3 & 0 & 0 & 1 & 1 & 0 \\ \hline
\end{osgridtabular}

结果是 \texttt{1000}。把任一进位线断开后重算；最先与正确答案不同的 bit
就是故障从哪一级开始传播。

\mustunderstand\textbf{第三步：用 NOR 门观察“保持”。}
交叉反馈的 SR 锁存器不依赖时钟：

\begin{osgridtabular}{cccl}
\(S\) & \(R\) & \(Q_{next}\) & 动作 \\ \hline
0 & 0 & \(Q_{old}\) & 保持以前的状态 \\ \hline
1 & 0 & 1 & set \\ \hline
0 & 1 & 0 & reset \\ \hline
1 & 1 & 未定义 & 两个输出不再互补 \\ \hline
\end{osgridtabular}

先施加 \(S=1,R=0\)，再回到 \(0,0\)。如果 Q 也跟着回到 0，反馈没有接好；
正确电路会继续保存 1。随后故意施加 \(1,1\)，记录两个输出和撤销顺序。这个
实验说明“未定义”不是随机挑一个正确答案。

\mustunderstand\textbf{第四步：把任意输入收束成 D 与 enable。}
令
\[
S=E\cdot D,\qquad R=E\cdot\overline D.
\]
当 \(E=1\) 时，Q 跟随 D；当 \(E=0\) 时，S=R=0，锁存器保持。因为 D 和
\(\overline D\) 不会同时为 1，这层组合逻辑避免正常使用时进入 SR 的非法
输入。让 E 保持 1 并翻转 D，会看到 Q 在整个窗口内透明；这还不是边沿触发。

\mustunderstand\textbf{第五步：只在上升沿更新。}
把两个互补相位的 latch 串起来，或在仿真器中用门级 D 触发器。按下面的输入
序列逐拍记录，不能在两次边沿之间偷改 Q：

\begin{osgridtabular}{L{0.16\linewidth}L{0.20\linewidth}L{0.20\linewidth}L{0.32\linewidth}}
时刻 & CLK 动作 & D & Q 应怎样变化 \\ \hline
t0 & 低电平 & 0 & 保持初值 0 \\ \hline
t1 & 上升沿前 & 1 & 尚未采样，Q 仍为 0 \\ \hline
t2 & 上升沿 & 1 & 经过 clock-to-Q 后变为 1 \\ \hline
t3 & 高电平中间 & 0 & Q 仍为 1 \\ \hline
t4 & 下一上升沿 & 0 & Q 变为 0 \\ \hline
\end{osgridtabular}

\mustunderstand\textbf{第六步：四只触发器组成寄存器。}
四只触发器共用 CLK 和 reset。每个 D 前放一个 MUX：
\[
D_i=
\begin{cases}
input_i,& load=1,\\
Q_i,& load=0.
\end{cases}
\]
先在一个边沿以 \(load=1,input=\texttt{1010}\) 写入，再令 load=0 并把
input 改成 \texttt{0101}。下一个边沿后 Q 仍应是 \texttt{1010}；若变成输入，
说明反馈候选没有接到 MUX。

\mustunderstand\textbf{第七步：把寄存器输出接回加法器。}
加法器另一输入固定为 \texttt{0001}，MUX 在 \(enable=1\) 时选择
\(Q+1\)，否则选择 Q。复位以后逐拍状态为：

\begin{osgridtabular}{L{0.10\linewidth}L{0.12\linewidth}L{0.17\linewidth}L{0.17\linewidth}L{0.17\linewidth}}
边沿 & enable & Q 边沿前 & D 候选 & Q 边沿后 \\ \hline
0 & -- & 未知 & reset & 0000 \\ \hline
1 & 1 & 0000 & 0001 & 0001 \\ \hline
2 & 1 & 0001 & 0010 & 0010 \\ \hline
3 & 0 & 0010 & 0010 & 0010 \\ \hline
4 & 1 & 0010 & 0011 & 0011 \\ \hline
5 & 1 & 0011 & 0100 & 0100 \\ \hline
\end{osgridtabular}

\begin{workedexample}
从 \texttt{1110} 连续计数两拍。第一拍加法器得到 \texttt{1111}，
外部进位为 0；第二拍低四位得到 \texttt{0000}，外部进位为 1。若计数器只
保存四位，Q 回绕到 0；如果还要检测溢出，就把外部进位保存到单独状态位。
\end{workedexample}

\begin{faultinjectionbox}{让状态链在不同位置断一次}
先把半加器的 XOR 错换成 OR，输入 \(1+1\) 会得到和位 1、进位 1，即错误值
\texttt{3}。恢复以后，再断开寄存器 MUX 的 Q 反馈：enable=0 时 D 不再有
确定来源，下一个边沿后的 Q 不能继续标为旧值。最后把 counter 的 enable
选择反接，比较上表第 3、4 拍。每次只引入一个错误，并记录第一个错误节点，
不要只记录最终计数不对。
\end{faultinjectionbox}

\begin{experimentbox}{自己增加“向下计数”输入}
在加法器第二输入前再放一个 MUX：up=1 选择 \texttt{0001}，up=0 选择
\texttt{1111}，后者是四位的 \(-1\)。验证
\texttt{0010}\(\rightarrow\)\texttt{0001}\(\rightarrow\)\texttt{0000}
\(\rightarrow\)\texttt{1111}，并解释最后一步为何按模 16 回绕。完成后画出
PC 使用同类“保持、加一、跳转目标、返回地址”四选一结构。
\end{experimentbox}

\topic{复位、时钟和 ready 的顺序}

\begin{osgridlongtable}{L{0.17\linewidth}L{0.22\linewidth}L{0.20\linewidth}L{0.31\linewidth}}
阶段 & 电源/时钟 & reset & 允许动作 \\ \hline
\endfirsthead
阶段 & 电源/时钟 & reset & 允许动作 \\ \hline
\endhead
上电初期 & 电源爬升，时钟未锁定 & 有效 & 寄存器保持复位值 \\ \hline
时钟稳定 & 周期满足规格 & 仍有效 & 等待同步释放 \\ \hline
复位释放 & 稳定边沿到达 & 撤销 & PC 取得复位向量 \\ \hline
首条总线读 & 地址与 read 稳定 & 无效 & 等待 ROM ready \\ \hline
ready 有效 & 数据满足采样 & 无效 & IR 接收，PC 才递增 \\ \hline
\end{osgridlongtable}

先释放 reset 再等待时钟，寄存器可能在不满足时序的边沿推进；ready 尚未有效
就递增 PC，返回 byte 会与错误地址配对。

\begin{experimentbox}{观察两拍 ready 延迟}
进入配套源码的 \filepath{examples/teaching_cpu} 目录，运行
\code{python3 verify.py}，找到 \code{scenario=ready-delay}。任取一个
fetch 地址，应先出现两行
\code{ready=0}，第三行同一地址才出现 \code{ready=1}。等待期间 PC、地址和
read 不得变化；正常与延迟场景最终都应写出
\code{memory[0x8000]=0x34}。
\end{experimentbox}

\begin{faultinjectionbox}{删除上拉并提前释放复位}
删除输入上拉后，把节点状态标为“未知”，不能在真值表中擅自填 0。再把复位
释放移到时钟稳定之前，找出第一个可能推进的寄存器并标未知。PC 一旦未知，
后续取指地址都无法计算。恢复上拉和释放顺序后，PC 才能从复位向量推进。
\end{faultinjectionbox}

\begin{pitfallbox}
\begin{itemize}[leftmargin=2.2em]
  \item 电压必须写清测量的两个点；GND 是选定参考网络；
  \item 上拉只在没有更强驱动时把节点带入高电平范围；
  \item 真值表没有表达传播延迟；
  \item latch 在窗口内透明，D 触发器在边沿附近采样；
  \item reset 定初值，clock 定推进时刻，ready 定事务完成时刻。
\end{itemize}
\end{pitfallbox}

\topic{练习：从引脚算到 PC}

\begin{enumerate}[leftmargin=2.2em]
  \item \(\SI{5}{V}\) 配 \(\SI{4.7}{k\ohm}\) 上拉，按键闭合时电流和功耗
    约为多少？
  \item \(R=\SI{22}{k\ohm}\)、\(C=\SI{100}{nF}\)，求 RC 和充到
    \(0.7V_{CC}\) 的时间。
  \item 求全加器 \(A=1,B=0,C_{in}=1\) 的两个输出。
  \item 边沿在 \(\SI{12}{ns}\)，setup 为 \(\SI{0.8}{ns}\)，hold 为
    \(\SI{0.3}{ns}\)。D 在什么区间必须稳定？
  \item PC MUX 的顺序值为 \texttt{1014}，目标为 \texttt{2000}。写出
    Z=0 和 Z=1 时的下一 PC。
\end{enumerate}

\begin{answerbox}
\begin{enumerate}[leftmargin=2.2em]
  \item \(I\approx\SI{1.064}{mA}\)，\(P\approx\SI{5.32}{mW}\)；
  \item \(RC=\SI{2.2}{ms}\)，70\% 约需 \(1.204RC=\SI{2.65}{ms}\)；
  \item \(S=0,C_{out}=1\)；
  \item \(\SI{11.2}{ns}\) 到 \(\SI{12.3}{ns}\)；
  \item Z=0 选 \texttt{1014}，Z=1 选 \texttt{2000}。
\end{enumerate}
\end{answerbox}
